\documentclass{article}
\usepackage{amsmath}
\usepackage{amssymb}

\begin{document}

\section*{Canonical Routing}

With canonical routing,
\begin{equation*}
p_i(b) = G_i(b),
\end{equation*}
agent \(i\) receives samples according to
\begin{equation*}
B(b)\,G_i(b).
\end{equation*}

In \texttt{position\_only}, the centroid weight is
\begin{equation*}
w(b) = 1 - G_i(b).
\end{equation*}
Therefore the target centroid is
\begin{equation*}
\hat{x}_i
=
\frac{
\sum_b B(b)\,G_i(b)\bigl(1-G_i(b)\bigr)\,b
}{
\sum_b B(b)\,G_i(b)\bigl(1-G_i(b)\bigr)
}.
\end{equation*}

The update rule is
\begin{equation*}
x_i^{+} = (1-\beta)x_i + \beta \hat{x}_i.
\end{equation*}
Hence
\begin{equation*}
\Delta x_i
=
x_i^{+}-x_i
=
\beta\,
\frac{
\sum_b B(b)\,G_i(b)\bigl(1-G_i(b)\bigr)\,(b-x_i)
}{
\sum_b B(b)\,G_i(b)\bigl(1-G_i(b)\bigr)
}.
\end{equation*}

This matches the gradient direction up to:
\begin{enumerate}
\item left multiplication by the \(\Sigma\)-metric (since the exact gradient contains \(\Sigma^{-1}\)),
\item a positive scalar normalization
\begin{equation*}
\sum_b B(b)\,G_i(b)\bigl(1-G_i(b)\bigr),
\end{equation*}
\item step size \(\beta\).
\end{enumerate}

So for isotropic \(\Sigma\), this is essentially a preconditioned/scaled gradient step, which is why the update is called \emph{gradient-matched in direction}.

\end{document}